% !TEX root = ../main.tex
\section{Caso de prueba}

Se evalúa el siguiente estudiante:

\begin{longtable}{lr}
\toprule
Entrada & Valor\\
\midrule
Promedio actual & 58\\
Asistencia & 64\\
Entregas realizadas & 55\\
Participación & 48\\
Exámenes recientes & 52\\
\bottomrule
\end{longtable}

\subsection{Fuzzificación}

Con los parámetros definidos, se obtienen grados aproximados:

\begin{longtable}{lll}
\toprule
Variable & Conjunto & Grado\\
\midrule
Promedio & bajo & $0.133$\\
Promedio & regular & $0.650$\\
Asistencia & baja & $0.400$\\
Asistencia & media & $0.267$\\
Entregas & insuficientes & $0.333$\\
Entregas & parciales & $0.250$\\
Participación & baja & $0.133$\\
Participación & media & $0.400$\\
Exámenes & deficientes & $0.533$\\
Exámenes & regulares & $0.111$\\
\bottomrule
\end{longtable}

\subsection{Activación de reglas}

Ejemplos de activación:

\[
\alpha_{R1}=\min(0.133,0.400)=0.133
\]

\[
\alpha_{R2}=\min(0.133,0.333)=0.133
\]

\[
\alpha_{R4}=\min(0.267,0.650)=0.267
\]

\[
\alpha_{R7}=\min(0.650,0.111)=0.111
\]

Las reglas producen salidas recortadas en los conjuntos crítico, alto y medio. La agregación final se calcula por máximo punto a punto.

\subsection{Resultado}

Con discretización de paso 1 en el universo $[0,100]$, el motor calcula:

\[
y^*\approx 60.37
\]

El valor pertenece principalmente al conjunto de riesgo alto. Por tanto, el resultado final es:

\[
\text{Riesgo académico} = 60.37,\quad \text{etiqueta}=\text{alto}
\]

\section{Resultados}

El sistema permite observar cada etapa del razonamiento. Las reglas no se ocultan y el valor final puede rastrearse desde los grados de pertenencia hasta la salida agregada. Esto cumple el objetivo de construir un modelo transparente, determinista y defendible académicamente.
